% S5 proofs for Section 6
\section{Proofs for Section 6: inference}\label{sec:proofs-inference}

Throughout, (I1)--(I2) of the main text hold: conditional on the design (and the
district effects), the unit residuals are independent with mean zero, and the design
is fixed. The regression coefficients $\hat\beta$ are OLS for a fixed-dimensional
parameter with a full-rank design, so $\sqrt n(\hat\beta-\beta_0)$ is asymptotically
normal with a consistent HC covariance estimate (outcomes are bounded). Constraint terms
are $s(\beta)=S\beta$ (stacked), the target is $\varphi_0(\beta)+\omega^\top k$, and
$\theta_\lambda(\beta)=g_\lambda^\top\beta$ with $g_\lambda=\nabla\varphi_0+S^\top\lambda$.

\begin{proofof}{Lemma~\ref{P-lem:minlin}}
The LP $\max\{\omega^\top k:s+Ak\ge0\}$ has Lagrangian $\omega^\top k+\lambda^\top(s+Ak)$;
stationarity in the free $k$ requires $\omega+A^\top\lambda=0$, so the dual is
$\min\{\lambda^\top s:\lambda\ge0,\ A^\top\lambda=-\omega\}$, and feasibility with bounded
value gives equal optimal values. The dual feasible set lies in $\{\lambda\ge0\}$, so it
contains no line and is pointed; a linear objective bounded below attains its minimum
at one of its finitely many vertices. The dual feasible set does not depend on $\beta$,
because $A$ and $\omega$ are fixed by the design. So
$\varphi_U(\beta)-\varphi_0(\beta)=\min_{\lambda\in\Lambda}\lambda^\top s(\beta)$ wherever
the primal is feasible and bounded. A minimum of finitely many linear functions is
concave and piecewise linear. The lower endpoint follows by replacing $\omega$ with
$-\omega$.
\end{proofof}

\begin{proofof}{Proposition~\ref{P-prop:consistency}}
(c) A vertex of $\hat A_d$ is an observed composition, hence in $A_d$, so
$\hat A_d\subseteq A_d$, the constraints over $\hat A_d$ are a subset of those over $A_d$,
and the feasible set is weakly larger.

(a) By Lemma~\ref{P-lem:minlin},
$\varphi_U(\beta)-\varphi_U(\beta')=\varphi_0(\beta-\beta')+[\min_\lambda\lambda^\top
s(\beta)-\min_\lambda\lambda^\top s(\beta')]$, and the bracket is bounded by
$\max\|S^\top\lambda\|\,\|\beta-\beta'\|$ over the optimal duals at $\beta$ or $\beta'$.
For $\beta$ near $\beta_0$, the Slater point $\tilde k$ stays $\varepsilon/2$-strictly
feasible, and under strict feasibility the optimal duals are bounded:
$\|\lambda\|_1\le(\text{value}-\text{value at the strictly feasible point})/(\varepsilon/2)$.
Hence $|\hat\varphi_U-\varphi_U(\beta_0;\hat A)|\le L\|\hat\beta-\beta_0\|=O_p(n^{-1/2})$,
with $L$ uniform in the hull because the constraint set over $\hat A$ is a subset of that
over $A$.

(b) The feasible sets $\Theta(\hat A)$ are compact and convex, contain $\Theta(A)$ by (c),
and decrease as $\hat A$ grows. If $k\in\Theta(\hat A_p)$ for all $p$ and $v\in A$, take
$v_p\in\hat A_p$ with $v_p\to v$; the constraints are continuous in $v$, so $k$ satisfies
the constraint at $v$, and $\bigcap_p\Theta(\hat A_p)=\Theta(A)$. A decreasing sequence of
compact convex sets converges in Hausdorff distance to its intersection, so the support
functions, and hence the LP values, converge. In the non-nested case,
$\Theta(A)\subseteq\Theta(\hat A_n)$ for every $n$; if $k_n\in\Theta(\hat A_n)$ and
$k_{n_j}\to k$, then for each $v\in A$ pick $v_j\in\hat A_{n_j}$ with $v_j\to v$, and
continuity of the constraints in $(k,v)$ gives $k\in\Theta(A)$. So the upper Kuratowski
limit is $\Theta(A)$, which is contained in every term; in a common bounded set this is
Hausdorff convergence.
\end{proofof}

The error splits as $\hat\varphi_U-\varphi_U(\beta_0;A)=[\hat\varphi_U-\varphi_U(\beta_0;\hat
A)]+[\varphi_U(\beta_0;\hat A)-\varphi_U(\beta_0;A)]$. The first part is $O_p(n^{-1/2})$ by
(a); the second is non-negative and tends to zero by (b), but no rate is shown for it.

\begin{proofof}{Theorem~\ref{P-thm:limit}}
By Lemma~\ref{P-lem:minlin}, $\varphi_U(\beta)=\min_{\lambda\in\Lambda}g_\lambda^\top\beta$,
a finite minimum of linear maps. Such a map is Lipschitz, and its directional derivative
at $\beta_0$ in direction $h$ is $\min_{\lambda\in\Lambda^*}g_\lambda^\top h$; a Lipschitz
map with directional derivatives in finite dimensions is Hadamard directionally
differentiable. The directional delta method gives the limit. (i) A minimum over a
singleton is linear, so the limit is normal. (ii) Let $Y=g_\lambda^\top\mathcal Z$ and
$Y'=g_{\lambda'}^\top\mathcal Z$ with $\Var(Y-Y')>0$. Then
$\E\min_{\Lambda^*}g^\top\mathcal Z\le\E\min(Y,Y')=\tfrac12\E(Y+Y')-\tfrac12\E|Y-Y'|=-\tfrac12\E|Y-Y'|<0$.
If $\Var(g_\lambda^\top\mathcal Z-g_{\lambda'}^\top\mathcal Z)=0$ for all active pairs, the
minimum equals each $g_\lambda^\top\mathcal Z$ a.s., which is normal. (iii) In case (ii)
the derivative is not linear on the support of $\mathcal Z$, so the condition of
Fang and Santos for bootstrap consistency fails. This theorem applies published results
to the structure of Lemma~\ref{P-lem:minlin}; its status is ``cited structure''.

The condition that the constraint system be fixed matters. For
$f_n(\theta)=\min\{\theta,-\theta+n^{-1/2}\}$ at $\theta_0=0$ with
$\sqrt n(\hat\theta-\theta_0)\Rightarrow Z$, the first piece is uniquely active for every
$n$, yet $\sqrt n(f_n(\hat\theta)-f_n(\theta_0))\Rightarrow\min\{Z,1-Z\}$, which is not
normal; Proposition~\ref{P-prop:triangular} covers this case.
\end{proofof}

\begin{proofof}{Proposition~\ref{P-prop:triangular}}
By Lemma~\ref{P-lem:minlin} for the design-$n$ system,
$\varphi_{U,n}(\beta)=\min_{\lambda\in\Lambda_n}\theta_{\lambda,n}(\beta)$. Hence, exactly,
$\sqrt n(\hat\varphi_{U,n}-\varphi_{U,n}(\beta_0))=\min_\lambda\big(g_{\lambda,n}^\top\sqrt
n(\hat\beta-\beta_0)+c_{\lambda,n}\big)$. Since $\sqrt n(\hat\beta-\beta_0)=O_p(1)$ and the
gradients are bounded, a candidate with $c_{\lambda,n}\to\infty$ exceeds the active
candidates' terms with probability tending to one and drops out. The remaining
finitely many terms converge jointly in distribution to
$(g_\lambda^\top\mathcal Z+c_\lambda)_\lambda$, and the minimum is continuous. Without the
localisation condition the same holds along every subsequence on which it holds, and
the limit may differ between subsequences.
\end{proofof}

\begin{proofof}{Theorem~\ref{P-thm:ci}(a)}
(1) Let $\mathcal E$ be the event that $s_j(\beta_0)\le s_j(\hat\beta)+\hat
c\hat\sigma_{s_j}$ for all $j$ and $|\varphi_0(\beta_0)-\varphi_0(\hat\beta)|\le\hat
c\hat\sigma_{\varphi_0}$. These are linear functionals of one asymptotically normal
vector, $\mathcal E$ is a convex event, the CLT holds uniformly over convex sets, and
$\hat\Sigma$ is consistent, so $\liminf P(\mathcal E)\ge1-\alpha_U-\alpha_L$. For M1 and
M2 this needs $\dim\beta$ fixed (regime R1).
(2) On $\mathcal E$, the true $k_0$ satisfies $s(\beta_0)+Ak_0\ge0$, hence
$s(\hat\beta)+\hat c\hat\sigma_s+Ak_0\ge0$, so $k_0\in\hat\Theta^+$. This uses only that
$k_0$ is feasible, not Slater's condition.
(3) On $\mathcal E$, $\Bt=\varphi_0(\beta_0)+\omega^\top k_0$ lies in
$[\varphi_0(\hat\beta)-\hat c\hat\sigma_{\varphi_0}+\min_{\hat\Theta^+}\omega^\top k,\
\varphi_0(\hat\beta)+\hat c\hat\sigma_{\varphi_0}+\max_{\hat\Theta^+}\omega^\top k]$.
(4) $E=\sum_iw_ie_i$ with $e_i$ independent given the design, mean zero and
$\Var(e_i\mid V_i)\le\tfrac14$ (Popoviciu), so $s_n^2:=\Var(E\mid\text{design})\le S_1^2$.
(5) $\limsup_nP(|E|>z_{\alpha_E}S_1)\le\alpha_E$ if $\max_iw_i^2/\sum_iw_i^2\to0$. It
suffices that every subsequence has a further subsequence along which this holds; take
one along which $s_n/S_1\to\rho\in[0,1]$. If $\rho=0$, Chebyshev's inequality gives
$P(|E|>z_{\alpha_E}S_1)\le s_n^2/(z_{\alpha_E}^2S_1^2)\to0$. If $\rho>0$, the summands
$w_ie_i$ are independent with $|w_ie_i|\le|w_i|$, and
$\max_i|w_i|/s_n\le2\max_i|w_i|/(\rho'(\sum_iw_i^2)^{1/2})\to0$ for any $0<\rho'<\rho$ and
large $n$; so Lindeberg's condition holds relative to $s_n$, $E/s_n\Rightarrow N(0,1)$,
and $P(|E|>z_{\alpha_E}S_1)\le P(|E|/s_n>z_{\alpha_E})\to\alpha_E$. (A central limit
theorem under the weight condition alone would not suffice, for instance when only one
summand has positive variance; the subsequence argument covers that case.)
(6) Since $B=\Bt+E$, $B\notin\CI$ requires either $\Bt$ outside the interval of (3) or
$|E|>z_{\alpha_E}S_1$. (7) A union bound gives
$\limsup P(B\notin\CI)\le\alpha_U+\alpha_L+\alpha_E$.
\end{proofof}

\emph{Part (b1)} replaces step (1) by CLR's coverage of each intersection bound, which
needs the finite index set $\Lambda$ of Lemma~\ref{P-lem:minlin} and their Condition~P;
steps (4)--(7) are unchanged. Condition~P(iv) requires gradients bounded away from zero
and a positive definite limiting covariance: in this setting, $C-1$ columns, a full-rank
design and $g_\lambda\ne0$ for every candidate.

\begin{proofof}{Theorem~\ref{P-thm:ci}(b2)}
(1) Under the model the primal LP at $\beta_0$ is feasible, because $k_0$ is. For any
dual-feasible $\lambda$ and any primal-feasible $k$,
$\omega^\top k=-\lambda^\top Ak\le\lambda^\top s(\beta_0)$, since $\lambda\ge0$ and
$s+Ak\ge0$. So $\theta_\lambda(\beta_0)\ge\varphi_U(\beta_0)$ for every dual-feasible
$\lambda$, not only optimal ones.
(2) On $\mathcal E=\{\theta_\lambda(\beta_0)\le\hat\theta_\lambda+\hat c\hat\sigma_\lambda\
\forall\lambda\in\bar\Lambda\}$,
$\hat U\ge\min_{\lambda\in\hat\Lambda}\theta_\lambda(\beta_0)\ge\varphi_U(\beta_0)$, because
$\hat\Lambda\subseteq\bar\Lambda$.
(3) Case (i): the standardised errors $(\theta_\lambda(\beta_0)-\hat\theta_\lambda)/\hat\sigma_\lambda$,
$\lambda\in\bar\Lambda$, are a fixed finite collection of standardised linear
functionals of $\hat\beta$, jointly asymptotically normal, so $P(\mathcal E)\to1-\alpha_U$.
(4) Case (ii): the reduced parameterisation contains the positive-definite case
($L=I$, $b=0$). Write $\hat e=\hat\eta-\eta_0$. For every $g$,
$g^\top(\hat\beta-\beta_0)=(L^\top g)^\top\hat e$ exactly, so by Cauchy--Schwarz in the
metric of $\hat\Sigma_\eta$ (positive definite with probability tending to one),
\[
|g^\top(\hat\beta-\beta_0)|\le(\hat e^\top\hat\Sigma_\eta^{-1}\hat e)^{1/2}(g^\top L\hat\Sigma_\eta L^\top g)^{1/2}
=(n\,\hat e^\top\hat\Sigma_\eta^{-1}\hat e)^{1/2}\hat\sigma_g .
\]
This holds for every $g$ at once, including $g$ with $L^\top g=0$, for which both sides
are zero. By continuous mapping $n\hat e^\top\hat\Sigma_\eta^{-1}\hat e\Rightarrow\chi^2_q$,
because $\Sigma_\eta$ is non-singular. So the event holds for every candidate and both
ends with probability tending to $1-\alpha$. Non-singularity is essential: with a
pseudoinverse of a singular $\Sigma$, an error of order $o(n^{-1/2})$ in a null direction
is neither zero nor controlled, as the example $(Z/\sqrt n,-1/n)$ of the main text shows.
\end{proofof}

\begin{proofof}{Lemma~\ref{P-lem:reduction}}
Write $\tau^*=\min_{\mu\in\Lambda}(\hat\theta_\mu+\hat c_p\hat\sigma_\mu)$ for the threshold
of the selection rule over $\Lambda$. A minimiser of this expression belongs to
$\hat\Lambda$, so the minimum over $\hat\Lambda$ equals $\tau^*$: the threshold is the same.
The set selected over $\hat\Lambda$ is CLR's selected set $\hat V$ intersected with
$\hat\Lambda$, and $\hat V\subseteq\hat\Lambda$ by hypothesis, so it is $\hat V$. The
critical value $\hat c$ is computed over the selected set, so it coincides, and so does
$\hat U=\min_{\lambda\in\hat V}(\hat\theta_\lambda+\hat c\hat\sigma_\lambda)$. The coverage of
Theorem~\ref{P-thm:ci}(b1) therefore carries over.
\end{proofof}

\begin{proofof}{Proposition~\ref{P-prop:bd}}
(1) $\bar m$ is a linear target under the rung (Proposition~\ref{P-prop:lp} with weights
$\pi$), and the rung alone holds at the truth, so $\mathcal E_v=\{\bar m\in\CI_v\}$ has
$\liminf P(\mathcal E_v)\ge1-\alpha_v$ by Theorem~\ref{P-thm:ci}(a) or (b2)(ii).
(2) By the Bhatia--Davis inequality and concavity of $m(1-m)$,
$s_n^2\le\sum_iw_i^2m_i(1-m_i)\le(\sum_iw_i^2)\bar m(1-\bar m)=\bar s^2$; on $\mathcal E_v$,
$\hat S_{BD}\ge\bar s\ge s_n$.
(3) So $\{|E|>z_{\alpha_E}\hat S_{BD}\}\subseteq\{|E|>z_{\alpha_E}\bar s\}\cup\mathcal E_v^c$, with
$\bar s$ deterministic given the design.
(4) $\limsup P(|E|>z_{\alpha_E}\bar s)\le\alpha_E$: along a subsequence with
$s_n/\bar s\to0$ use Chebyshev; otherwise $\max_i|w_i|/s_n\to0$ by the condition
$\max_iw_i^2/\bar s^2\to0$, and Lindeberg's theorem applies relative to $s_n$, as in step
(5) of part (a) with $\bar s$ in place of $S_1$.
(5) A union bound with the estimation event of (a) or (b2)(ii) finishes.
If $\CI_v$ is a CLR interval, the same argument holds on the event that $\CI_v$ covers
$\bar m$, whose probability is the subject of Conjecture~\ref{P-conj:clr}.
\end{proofof}

\begin{proofof}{Proposition~\ref{P-prop:tests}}
(a) Under the model $k_0\in\hat\Theta^+$ on $\mathcal E$ (step (2) of part (a) of
Theorem~\ref{P-thm:ci}), so $P(\hat\Theta^+=\emptyset)\le P(\mathcal E^c)$, whose limit
superior is at most $\alpha_U+\alpha_L$.

(b) This is the standard Wald test for OLS with a fixed number of restrictions; the
count of restrictions follows from Lemma~\ref{P-lem:D1}(b): the rank gap between the M2
and M1 designs is $(D^\star-1)R(R-1)/2$ per column, and column $C$'s restrictions are
implied by adding up.

(c) Under M1 and (I1), $\E[\beta_i\mid\text{design}]=M_H(X_i)+Z_{D_i}$ with
$Z_d=[\zeta_{c,d,r}]_{rc}$. M0 holds for the observed units iff some $H'$ has
$M_{H'}(X_i)=M_H(X_i)+Z_{D_i}$ for all $i$, i.e.\ $G_cX_i=\zeta_{c,D_i}$ with $G=H'-H$.
Within district $d$ this reads $(G_c-\zeta_{c,d}\one^\top)X_i=0$ for every unit $i$ in
$d$, using $\one^\top X_i=1$. The district's compositions include $R$ affinely
independent points, which are linearly independent because they lie in
$\{\one^\top x=1\}$, so $G_c=\zeta_{c,d}\one^\top$. This holds for every $d$ and
$\zeta_{c,1}=0$, so $G_c=0$ and every $\zeta_{c,d}=0$. Conversely $\zeta=0$ gives M0 with
$H'=H$. The shifts are coefficients of the M1 regression, identified by
Lemma~\ref{P-lem:D1}(d), and the Wald test of $\zeta=0$ has $(D^\star-1)(C-1)R$
restrictions for the $C-1$ free columns.
\end{proofof}

\begin{proofof}{Theorem~\ref{P-thm:finite}}
(1) $T_i\in\Delta^{C-1}$, so $Y_i(f)=\sum_cw_{ic}(f)T_{ic}$ is a convex combination of the
$w_{ic}(f)$ and lies in an interval of length $\rho_i(f)$.
(2) Under the mean model and full rank, OLS is conditionally unbiased:
$\E[f(\hat\beta)\mid\text{design}]=f(\beta_0)$, because $\E[T_{ic}\mid\text{design}]=z_i^\top\beta_{c,0}$.
(3) The $Y_i(f)$ are independent given the design, so by Hoeffding's inequality
$P(f(\beta_0)-f(\hat\beta)\ge h(f;\delta))\le\delta$, and likewise on the other side.
(4) Let $\mathcal E$ be the event that $s_j(\beta_0)\le s_j(\hat\beta)+h_j$ for all $j$ and
$|\varphi_0(\beta_0)-\varphi_0(\hat\beta)|\le h_0$. It is an intersection of $J+2$ one-sided
events, so $P(\mathcal E^c)\le(J+2)\delta=\alpha_s$.
(5) On $\mathcal E$, the true $k_0$ satisfies $s(\beta_0)+Ak_0\ge0$, hence
$k_0\in\hat\Theta^{\rm fin}$, and $\Bt=\varphi_0(\beta_0)+\omega^\top k_0$ lies in the interval
without $t_E$. That proves (b), and (c), because $\hat\Theta^{\rm fin}\ne\emptyset$ on
$\mathcal E$.
(6) $E=\sum_iw_{ir}e_{irc}$ with independent $e_{irc}\in[-m_{rc}(V_i),1-m_{rc}(V_i)]$, an
interval of length one, so $P(|E|\ge t_E)\le2\exp(-2t_E^2/\sum_iw_{ir}^2)=\alpha_E$.
(7) $B=\Bt+E$, and a union bound over (4) and (6) gives (a). (d) holds because $B$ lies
in the Duncan--Davis interval surely.
\end{proofof}

\begin{proofof}{Proposition~\ref{P-prop:select}}
(a) Fewer constraints give a weakly larger feasible set (Theorem~\ref{P-thm:ladder}(b)).
(b) Conditioning on the design fixes $\mathcal D'$, so a procedure's conditional coverage
given the design, for the fixed set $\mathcal D'$, is its coverage given the design and
$\mathcal D'$.
\end{proofof}

\begin{proofof}{Proposition~\ref{P-prop:S1}}
(1) $\E[T\mid X]=X^\top QX+X^\top\nu(X)=X^\top QX$, so $Q$ is unchanged.
(2) $M_H(x)+\nu(x)\in\Bset$ implies $M_H(x)\ge-\nu(x)\ge-\eta$ elementwise on $A$; $M_H$
is linear, so this is equivalent to the same inequality at the vertices,
$s+Ak\ge-\eta\one$.
(3) $\Bt_{rc}=\sum_iw_i[M_H(X_i)+\nu(X_i)]_{rc}=\varphi(k)+\sum_iw_i\nu_{rc}(X_i)$, and the last
term lies in $[-\eta,\eta]$ because the weights sum to one. For $\sum a_{rc}B_{rc}$ the
last term lies in $[-\eta\sum|a_{rc}|,\eta\sum|a_{rc}|]$.
The breakdown value is well defined because the bounds are nested in $\eta$: a
conclusion that holds for every value in the $\eta$-bound holds for every smaller $\eta$.
\end{proofof}

\subsection{Inference with added rows}

\begin{proofof}{Theorem~\ref{P-thm:R10}}
(a) (1) Under the model, $k_0$ satisfies the base rows and the added rows at $\beta_0$
(Theorem~\ref{P-thm:R3}(e) and the hypothesis). (2) Each added $s_j(\hat\beta)$ is a linear
map of $\hat\beta$, so it is jointly asymptotically normal with the base terms; with the
maximum in step (1) of Theorem~\ref{P-thm:ci}(a) taken over all rows,
$\liminf P(\mathcal E)\ge1-\alpha_U-\alpha_L$. A row whose $s_j$ does not depend on $\beta$
has zero variance, is imposed exactly, and is left out of the maximum. (3) On $\mathcal E$,
$k_0\in\hat\Theta^+$, and steps (3)--(7) of Theorem~\ref{P-thm:ci}(a) and
Proposition~\ref{P-prop:tests}(a) go through, since they use only $k_0\in\hat\Theta^+$ and
the residual term. (4) For version (b), the endpoint is still a minimum of finitely
many linear functions of $\beta$, because Lemma~\ref{P-lem:minlin} needs only rows linear
in $\beta$ with fixed $a_j$. CLR applies to the augmented system if its conditions hold
there: the added rows change the dual feasible set, its vertices, the gradients, the
selection and the quantile, so the conditions must be checked anew; the projection and
fixed-superset arguments of (b2) carry over directly. For a contrast $B_1-B_0$, the
residual is $\sum_i(w_{i1}e_{i1}-w_{i0}e_{i0})$, each unit term has range at most
$w_{i1}+w_{i0}$, so $\tfrac12\{\sum_i(w_{i1}+w_{i0})^2\}^{1/2}$ is a valid standard-deviation
envelope under conditional independence; the sum of the per-cell envelopes is also
valid, by Minkowski's inequality, and more conservative.

(b) (5) Let $l^*$ be a piece that holds at the truth, $s_{l^*}(\beta_0)+a^\top k_0\ge0$.
(6) On $\mathcal E$, which now includes every piece,
$s_{l^*}(\beta_0)\le s_{l^*}(\hat\beta)+\hat c\hat\sigma_{l^*}$, so
$\max_l(s_l(\hat\beta)+\hat c\hat\sigma_l)+a^\top k_0\ge0$. (7) For given data this is one
linear row in $k$, which $k_0$ satisfies on $\mathcal E$; proceed as in (3). The argument
uses one simultaneous event and no strict feasibility. At a tie, inflating all pieces
and then maximising still dominates the true piece. The row is not linear in $\beta$
(an endpoint becomes a maximum of minima), so Lemma~\ref{P-lem:minlin} fails and CLR does
not apply; for CLR the branch union of (c) is used. (8) For one-vs-rest and $2\times2$
CR, Theorem~\ref{P-thm:R3}(d) gives
$D^{r*}\ge\min(\ER^*,\gamma_r\ER^*)\iff\max(D^{r*}-\ER^*,D^{r*}-\gamma_r\ER^*)\ge0$, and
similarly for the upper side; each piece is linear in $\beta$ apart from the common
term in $k$.

(c) (9) Fix a branch $j^*$ with $k_0\in C_{j^*}$; it is non-random, because $k_0$ is a
population quantity. (10) The model ``rung $\cap\,C_{j^*}$'' holds, so by (a),
$\liminf P(\theta_0\in\CI(C_{j^*}))\ge1-\alpha$. (11)
$\CI(C_{j^*})\subseteq\CI(C_+)\cup\CI(C_-)\subseteq$ its hull. (12) The other branch's
interval is never used for coverage, so its level is irrelevant, and no split is
needed. (13) Both sets empty implies the $j^*$ set is empty, which has limiting
probability at most $\alpha_U+\alpha_L$ (Proposition~\ref{P-prop:tests}(a) for version (a),
the crossing argument of (d) for version (b)). If the true branch switches along a
triangular array of designs, branch validity is needed along each subsequence on which
a branch is true; there are only two labels.

(d) (14) For version (a), $\CI=\emptyset$ requires $\hat\Theta^+=\emptyset$, which requires
$\mathcal E^c$. (15) For version (b) with $\theta_0=\Bt$, a crossing $\hat L>\hat U$ requires
$\hat U<\theta_0$ or $\hat L>\theta_0$, with limiting probabilities at most $\alpha_U$ and
$\alpha_L$ by one-sided validity, whatever the width of the identified set. For the
target $B$, the residual allowance only widens the interval, so a crossing of the
$B$-interval implies a crossing of the $\Bt$-interval. (16) In both cases the empty event
lies inside a non-coverage event of bounded probability. (17) Since
$\{\theta_0\in\CI\}\subseteq\{\CI\ne\emptyset\}$,
$P(\text{cover}\mid\CI\ne\emptyset)=P(\text{cover})/P(\CI\ne\emptyset)\ge P(\text{cover})$. If the
restrictions are false, step (1) fails and nothing is claimed.

(e) (18) Repeat each of $X=1/4,1/2,3/4$ equally often with equal sizes. Take $m_0=.4$
and $m_1=.55$ constant, with group rates drawn independently across units and groups
from bounded laws with these means and positive variance (for example, beta draws of
fixed concentration). Every base row has strict slack, OLS has a full-rank fixed design
and a positive definite limiting covariance, and HC3 is consistent. The positive CR
branch has $\tilde D_0=\ER=.15$ on its upper edge. (19) In the basis
$((1-x)^2,2x(1-x),x^2)$, for any coefficient vector $q$, $\ER(q)=q_{11}-q_{00}$ and the
target constant is $c_D(q)=\tfrac7{12}(q_{11}-q_{00})$, because $\E[X^2]/\E[X]=7/12$ and the
middle coefficient cancels by symmetry; so the ratio of their standard errors is
exactly $7/12$. (20) With probability tending to one the plug-in sign is positive and
the base rows allow the $k$ attaining the plug-in upper edge (strict slack and
continuity), so the naive upper endpoint is $\widehat\ER+k_{\rm crit}\operatorname{se}(c_D)$.
(21) Its critical value is calibrated on eight one-sided base rows and two target tails,
so by the union bound $k_{\rm crit}\le z_{.995}$ at level .95, whatever the correlations.
(22) Non-coverage is the event $\{\widehat\ER-\ER<-k_{\rm crit}\operatorname{se}(c_D)\}$, so
$\limsup_nP(\tilde D_0\in\CI_{\rm naive})\le\Phi(\tfrac7{12}z_{.995})=0.933525<0.95$.
(23) With a simulated critical value from 4000 draws, each draw exceeds $z_{.996}$ with
probability at most .04 by the same union bound, and an empirical .95 quantile above that
threshold needs at least 200 exceedances; so the limiting coverage is at most
$\Phi(\tfrac7{12}z_{.996})+P\{\mathrm{Bin}(4000,.04)\ge200\}=0.940092$.

For the literal intersection of the fitted bracket with a finished JKST interval: in the
design of (e), its upper endpoint cannot exceed $\widehat\ER$ with probability tending to
one, so its limiting coverage of $\tilde D_0=\ER$ is at most one half. The intersection
is not in general contained in the naive interval: dropping an empty intersection
before expanding endpoints and expanding first do not commute (with a raw base interval
$[0,.1]$, bracket $[.15,.2]$ and target allowance .1, the raw intersection is empty
although the expanded base meets the bracket), so this bound is proved directly.
\end{proofof}
